\section{System and calibration task}
\label{sec:system}

\subsection{Flux-tunable transmon and calibration objective}
\label{sec:system-objective}

We consider a flux-tunable transmon qubit~\cite{krantz2019quantum} controlled through a calibrated flux-bias coordinate $\Phi$. 
%Experimentally, $\Phi$ is controlled through a room-temperatue bias voltage $V_{\mathrm b}$ using a separately characterized voltage-to-flux conversion, as detailed in the Supplemental Material.
In the transmon regime, the flux dispersion of the qubit angular transition frequency is approximately given by~\cite{koch2007charge}
%
\begin{equation}
\begin{split}
\hbar\omega_{01}(\Phi)
\simeq{}&
\sqrt{8E_CE_J(\Phi)}-E_C, \\
% E_J(\Phi)
% ={}&
% E_{J,\Sigma}
% \sqrt{
% \cos^2\left(\pi\Phi/\Phi_0\right)
% +
% d^2\sin^2\left(\pi\Phi/\Phi_0\right)
% },
\label{eq:transmon-dispersion}
\end{split}
\end{equation}

Equation~\eqref{eq:transmon-dispersion} describes the flux dependence of the
transition frequency and thus provides the physical basis for frequency
calibration through flux-bias tuning. The calibration objective is to find a
flux bias $\Phi_\star$ such that
%
\begin{equation}
\left|\omega_{01}(\Phi_\star)-\omega_{\mathrm{tar}}\right|
\leq \epsilon_\omega,
\label{eq:calibration-success}
\end{equation}
%
where $\omega_{\mathrm{tar}}$ is the target angular transition frequency and
$\epsilon_\omega$ is the prescribed tolerance. This condition defines calibration success; the subsequent workflow monitors
the verified working point and initiates independent verification or
recalibration when drift is suspected.

In practice, a previously characterized frequency--flux map may no longer be
accurate at the target tolerance because of slow shifts in the effective flux
offset or device parameters\cite{paladino2014noise,vepsalainen2022improving}. We therefore formulate the calibration without
assuming an up-to-date quantitative map $\omega_{01}(\Phi)$: the controller
operates in the calibrated flux coordinate and obtains the local frequency
information required for bias updates directly from measurements.

Figure~\ref{fig:system-protocol}(a) summarizes this calibration setting. The
full dispersion illustrates the physical tunability of the qubit, while the
inset shows a local trajectory from the seed bias toward the target frequency.
% The displayed dispersion is withheld from the controller, which uses only the
% frequency information obtained at the visited bias points.

\begin{figure*}[t]
    \centering
    \includegraphics[width=0.9\textwidth]{figures/fig1/fig1-system-protocol.pdf}
    \caption{Schematic overview of the physical system, measurement probes,
    and calibration protocol. (a) Qubit frequency dispersion, with an inset showing the seed bias, target
    frequency, and successive points of the local calibration trajectory. The
    complete dispersion is withheld from the calibration controller. (b) Timing diagrams for scanned Ramsey,
    fixed-delay Ramsey, and the bipartite short-pulse probe. (c) The
    \textsc{Acquire--Track--Verify--Lock} protocol, with
    \textsc{Reacquire} recovery and a terminal safe hold. Only
    \textsc{Track} changes the flux bias; \textsc{Verify} freezes the
    candidate bias for independent confirmation, and \textsc{Lock} performs
    long-term monitoring. Solid, dashed, and dotted arrows denote forward,
    fallback, and failure transitions, respectively.}
    \label{fig:system-protocol}
\end{figure*}

% A calibration cycle combines three distinct operations: measurement, estimation, and tuning. A measurement executes a probe at fixed $\Phi_n$ and returns population data. An estimator maps these data to $\widehat{\Delta}_n$ and $\widehat{\omega}_{01,n}$, and a tuning step changes the flux bias. Repeated cycles locate the target working point. Stabilization, in contrast, suppresses subsequent drift around an established working point and is not the primary task considered here.

\subsection{Ramsey and bipartite short-pulse probes}
\label{sec:system-probes}

At each calibration iteration, the qubit is interrogated at a fixed flux bias
by a control sequence followed by population readout. We refer to this
pulse-and-readout operation as a measurement probe. We consider two types of
probes: Ramsey probes and bipartite short-pulse probes. The driven dynamics of
both probes are described by an effective two-level Hamiltonian. In a frame
rotating at the drive frequency $\omega_d$ and under the rotating-wave
approximation~\cite{krantz2019quantum,blais2021circuit},
\begin{equation}
\begin{aligned}
\frac{H(t,\Phi)}{\hbar}
={}&
\frac{\Delta(\Phi)}{2}\sigma_z
+
\frac{\Omega(t)}{2}
\left[
\cos\varphi(t)\sigma_x
+
\sin\varphi(t)\sigma_y
\right],
\\
\Delta(\Phi)
={}&
\omega_{01}(\Phi)-\omega_d .
\end{aligned}
\label{eq:hamiltonian}
\end{equation}
%
Here $\Delta(\Phi)$ is the angular-frequency detuning, while $\Omega(t)$ and
$\varphi(t)$ are the drive envelope and phase. The detuning term connects the
unknown transition frequency to the driven dynamics. Ramsey and short-pulse
probes use the distinct control sequences shown in
Fig.~\ref{fig:system-protocol}(b) to infer this detuning from population
readout. once $\widehat\Delta$ is obtained, the transition frequency follows
as $\widehat\omega_{01}=\omega_d+\widehat\Delta$.

% Figure~\ref{fig:system-protocol}(b) compares the timing structures of the
% two probe types. We first consider the Ramsey probe, which uses two nominal
% $\pi/2$ pulses separated by a free-evolution interval
% $\tau$~\cite{ramsey1950molecular}.

We first consider the Ramsey probe, which uses two nominal $\pi/2$ pulses
separated by a free-evolution interval $\tau$~\cite{ramsey1950molecular}. For an ideal Ramsey sequence with analysis phase $\varphi$ and Ramsey-fringe visibility $C$, the measured population can be written as
\begin{equation}
    p_R(\Delta; \tau, \varphi) = \frac{1}{2}\left[ 1 + C\cos(\Delta\tau + \varphi) \right],
    \label{eq:ramsey-response}
\end{equation}
We use scanned Ramsey, which samples
multiple values of $\tau$, for wide-range frequency estimation. As a local
baseline, we instead fix $\tau=\tau_0$ and use two orthogonal analysis phases
to recover the detuning from measured quadratures within the unambiguous phase
range~\cite{vepsalainen2022improving}. 
% The data processing and choice of
% $\tau_0$ are specified in Appendix~\ref{app:methods}.

We next turn to the bipartite short-pulse probe. Unlike the Ramsey probe,
it contains no separate free-evolution interval and instead applies a common
preparation rotation of angle $\alpha$ immediately
followed by one of two analysis rotations whose phases differ by $\pi$~\cite{herb2024quantum}:
\begin{equation}
    R_y(\alpha) \longrightarrow R_{+x}(\alpha) \quad\text{or}\quad R_y(\alpha) \longrightarrow R_{-x}(\alpha).
    \label{eq:bipartite-sequences}
\end{equation}

Here $\alpha = \int \Omega(t) \dd t$ is determined from the pulse envelope and duration. Varying $\alpha$ changes the differential detuning response and therefore the response coefficients and validated operating range introduced in Sec.~III.

We refer to these two sequences as the $+X$ and $-X$ branches and denote
their final excited-state populations by $p_{+X}$ and $p_{-X}$,
respectively. We define the corresponding differential observable as
\begin{equation}
    p_{\mathrm d}
    =
    \frac{p_{+X}-p_{-X}}{2}.
    \label{eq:pdiff}
\end{equation}

Detuning acts during the finite control pulses and changes the measured branch
populations. Their difference provides a first-order-sensitive local frequency
observable when
\begin{equation}
    \left.
    \frac{\partial p_{\mathrm d}}{\partial\Delta}
    \right|_{\Delta=0}
    \neq 0.
    \label{eq:local-identifiability}
\end{equation}
This condition is necessary for local inversion but does not guarantee an
unambiguous response over an unrestricted detuning range. Section~III
quantifies the local response and establishes its validated operating interval.


\subsection{Calibration workflow}
%  and independent validation}
\label{sec:calibration-workflow}

The local-identifiability condition applies at a fixed bias, whereas a bias
update can move the qubit outside the validated short-pulse range. We therefore
embed the local probe in the
\textsc{Acquire--Track--Verify--Lock} workflow shown in
Fig.~\ref{fig:system-protocol}(c), supplemented by
\textsc{Reacquire} and terminal \textsc{SafeStop} states. This architecture is
motivated by automated qubit calibration and closed-loop frequency
feedback~\cite{kelly2018physical,wittler_integrated_2021,
vepsalainen2022improving}.

At iteration $n$, the measurement detuning and target residual are
%
\begin{equation}
\Delta_n=\omega_{01}(\Phi_n)-\omega_{d,n},
\qquad
r_n=\omega_{01}(\Phi_n)-\omega_{\mathrm{tar}}.
\label{eq:detuning-and-residual}
\end{equation}
%
The detuning $\Delta_n$ determines whether the local probe remains valid,
whereas $r_n$ determines whether the target has been reached. A flux-bias
update changes the physical working point to reduce $r_n$; a drive-reference
update controls the detuning of the next measurement without changing the
target.

For both \textsc{Acquire} and \textsc{Track}, an estimate becomes a target
candidate when
%
\begin{equation}
|\widehat r_{x,k}|
+z_{1-\beta}\sigma_{x,k}
\leq\epsilon_{\mathrm{enter}},
\qquad
x\in\{\mathrm{acq},\mathrm{tra}\},
\label{eq:target-entry-condition}
\end{equation}
%
where $\sigma_{x,k}$ is the uncertainty of the estimated target residual and
$z_{1-\beta}$ is the prescribed confidence multiplier.

In \textsc{Acquire}, scanned Ramsey estimates the current qubit frequency and
initializes the drive. An estimate satisfying
Eq.~\eqref{eq:target-entry-condition} is routed directly to
\textsc{Verify}; otherwise, it seeds \textsc{Track}. Only \textsc{Track}
updates the flux bias, using the short-pulse frequency estimate while updating
the drive reference to keep subsequent measurements within the validated
range. When a \textsc{Track} estimate satisfies
Eq.~\eqref{eq:target-entry-condition}, the bias is frozen and passed to
\textsc{Verify}. A failed local-validity check instead enters
\textsc{Reacquire}.

\textsc{Verify} uses independent Ramsey measurements at the frozen bias.
Calibration is confirmed only after $N_{\mathrm{verify}}$ consecutive
measurements satisfy
%
\begin{equation}
|\widehat r_{\mathrm{ver},j}|
+z_{1-\beta}\sigma_{\mathrm{ver},j}
\leq\epsilon_{\mathrm{final}}.
\label{eq:verification-condition}
\end{equation}
%
A correctable verification error returns to \textsc{Track}, whereas a lost
frequency reference or an error outside the local recovery range enters
\textsc{Reacquire}.

After successful verification, \textsc{Lock} holds the verified bias and uses
periodic low-cost short-pulse measurements to monitor drift. Detected drift or
a scheduled audit returns the protocol to \textsc{Verify}. Before a monitor
residual is used for routing, its applied bias and drive are checked and its
probe detuning must satisfy the same local-validity condition in
Eq.~\eqref{eq:track-range-condition}. A large frequency jump, a failed
monitor-validity check, or loss of the local reference enters
\textsc{Reacquire}.
\textsc{Reacquire} repeats the wide-range acquisition and applies the same
routing rule as \textsc{Acquire}. Cancellation, budget exhaustion, or a safety
interlock enters \textsc{SafeStop}. Detailed thresholds and schedules are
specified in Appendix~\ref{app:methods}.
% The local-identifiability condition above concerns a fixed bias, whereas a
% calibration update can move the qubit beyond the invertible short-pulse range.
% We therefore use the \textsc{Acquire--Track--Verify--Lock} workflow in
% Fig.~\ref{fig:system-protocol}(c), with \textsc{Reacquire} recovery and a
% terminal \textsc{SafeStop}. \textsc{Acquire} establishes a wide-range absolute
% reference; \textsc{Track} alone performs calibration bias updates;
% \textsc{Verify} independently confirms a frozen candidate; and \textsc{Lock}
% monitors without correcting. \textsc{Reacquire} restores a lost reference,
% whereas \textsc{SafeStop} applies a safe bias hold. This separation follows
% the general setting of automated qubit calibration and closed-loop frequency
% feedback~\cite{kelly2018physical,wittler_integrated_2021,
% vepsalainen2022improving}.

% A probe returns population data, an estimator converts them into a frequency
% estimate and uncertainty, and the controller may update the physical bias. A
% separate drive-reference update changes only the reference of the next
% measurement. These operations are governed by
% %
% \begin{equation}
% \Delta_n
% =
% \omega_{01}(\Phi_n)-\omega_{d,n},
% \qquad
% r_n
% =
% \omega_{01}(\Phi_n)-\omega_{\mathrm{tar}}.
% \label{eq:detuning-and-residual}
% \end{equation}
% %
% The measurement detuning $\Delta_n$ guards local-probe validity, whereas the
% residual $r_n$ tests the fixed target. A bias update reduces $r_n$ by moving
% the working point; a drive update controls $\Delta_n$ without reducing
% $r_n$ or redefining the target.

% With $z_{1-\beta}$ denoting the prescribed confidence multiplier, define
% %
% \begin{align}
% U_{x,k}&=|\widehat r_{x,k}|+z_{1-\beta}\sigma_{x,k},
% &x&\in\{\mathrm{acq},\mathrm{tra},\mathrm{ver},\mathrm{loc}\},
% \label{eq:track-verify-statistics}\\
% \epsilon_{\mathrm{hold}}&<\epsilon_{\mathrm{final}}
% <\epsilon_{\mathrm{enter}}<\Delta_{\mathrm{val}}.
% \label{eq:threshold-hierarchy}
% \end{align}
% %
% Here $U_{x,k}$ is always a target-residual statistic. In contrast,
% $\Delta_{\mathrm{val}}$ bounds the probe detuning in
% Eq.~\eqref{eq:track-range-condition}; it must not be applied directly to
% $|\widehat r_{\mathrm{tra},k}|$. The ordering in
% Eq.~\eqref{eq:threshold-hierarchy} compares policy scales, not interchangeable
% observables.

% \textsc{Acquire} routes a reliable off-target estimate to \textsc{Track}. An
% acquired estimate or an accepted \textsc{Track} estimate instead becomes a
% target candidate when
% %
% \begin{equation}
% U_{x,k}\leq\epsilon_{\mathrm{enter}},
% \label{eq:target-entry-condition}
% \end{equation}
% %
% for $x\in\{\mathrm{acq},\mathrm{tra}\}$. The bias is then frozen and the
% protocol enters \textsc{Verify}; neither state can declare success. A valid
% off-target \textsc{Track} estimate instead produces the bounded update of
% Sec.~\ref{sec:calibration-protocol} and remains in \textsc{Track}. A failed
% range, branch-continuity, or sensitivity check enters \textsc{Reacquire}
% rather than extrapolating the local inverse.

% \textsc{Verify} uses measurements independent of the short-pulse estimator
% and its convergence history. Calibration is confirmed only after
% $N_{\mathrm{verify}}$ consecutive measurements at the frozen bias satisfy
% %
% \begin{equation}
% U_{\mathrm{ver},j}\leq\epsilon_{\mathrm{final}}.
% \label{eq:verification-condition}
% \end{equation}
% %
% A reliable verification miss returns to \textsc{Track}; an ambiguous or
% locally invalid verification result enters \textsc{Reacquire}. Measurements
% collected in \textsc{Track} never count toward $N_{\mathrm{verify}}$, even
% when they satisfy Eq.~\eqref{eq:target-entry-condition}.

% Successful verification enters \textsc{Lock}, which retains the verified bias
% and seeks $U_{\mathrm{loc}}\leq\epsilon_{\mathrm{hold}}$. A clear monitor
% result remains in \textsc{Lock}; repeated suspicion or a scheduled Ramsey
% audit enters \textsc{Verify}, while a large jump or lost reference enters
% \textsc{Reacquire}. The latter repeats wide-range acquisition and routes its
% seed by the same rule as \textsc{Acquire}. Budget exhaustion, cancellation,
% or an interlock instead enters terminal \textsc{SafeStop}.

% Response calibration, held-out response validation, closed-loop verification,
% and long-term monitoring use distinct data roles. Their thresholds, budgets,
% and schedules are given in Appendix~\ref{app:methods}. 
% % The present study
% % reports deterministic solver cost and analytic residuals, not hardware
% % $T_{\mathrm{hit}}$ or long-term locking data.

% Section~\ref{sec:inference-calibration} constructs only the estimator and the
% bias- and drive-reference updates used within \textsc{Track}; all state
% transitions and the definition of calibration success remain those above.
